\documentclass[aspectratio=169,10pt]{beamer}
\makeatletter\def\input@path{{../estilo/}{estilo/}}\makeatother
\usepackage{devmobile}
\modulo{09}{Firebase: Backend como Serviço}{Authentication, Cloud Firestore, regras de segurança, Storage, Cloud Messaging, Crashlytics e alternativas}

\begin{document}
\capa

\begin{frame}{Objetivos do módulo}
  \begin{conceito}[Ao final deste módulo você será capaz de]
    \begin{itemize}
      \item Explicar o modelo \emph{Backend as a Service} e os principais serviços do Firebase.
      \item Configurar um projeto Firebase em um app Android.
      \item Autenticar usuários com e-mail/senha e conta Google.
      \item Modelar, gravar e observar dados em tempo real no Cloud Firestore.
      \item Proteger os dados com regras de segurança e testá-las no emulador local.
      \item Enviar arquivos ao Storage, receber notificações push e monitorar falhas.
    \end{itemize}
  \end{conceito}
\end{frame}

\sumario

% =============================================================================
\section{Visão geral}

\begin{frame}{\emph{Backend as a Service}}
  \begin{columns}[c]
    \begin{column}{0.5\textwidth}
      \centering
      \begin{tikzpicture}[node distance=3mm]
        \node (cel) {\celular[0.45]{\tiny LabMonitor}};
        \node[caixaambarforte,minimum width=3.3cm,right=10mm of cel,yshift=10mm] (auth) {Authentication};
        \node[caixaescura,minimum width=3.3cm,below=of auth] (fs) {Cloud Firestore};
        \node[caixapetroleo,minimum width=3.3cm,below=of fs] (st) {Cloud Storage};
        \node[caixa,minimum width=3.3cm,below=of st] (fcm) {Cloud Messaging};
        \node[pontilhado,fit=(auth)(fs)(st)(fcm),label={[rotulo]above:\textbf{Firebase (Google Cloud)}}] (fb) {};
        \draw[seta,<->] (cel.east) -- (fb.west);
      \end{tikzpicture}
    \end{column}
    \begin{column}{0.48\textwidth}
      \small
      O \termo{Firebase} oferece, como serviços gerenciados, as peças que todo app conectado precisa:
      \begin{itemize}
        \item login de usuários;
        \item banco de dados em tempo real com cache offline;
        \item armazenamento de arquivos;
        \item notificações push;
        \item relatórios de falhas e métricas de uso.
      \end{itemize}
      \begin{dica}
        Ótimo para protótipos, TCCs e apps pequenos e médios: você se concentra no app e não em servidores.
      \end{dica}
    \end{column}
  \end{columns}
\end{frame}

\begin{frame}{Serviços e planos}
  \begin{columns}[T]
    \begin{column}{0.56\textwidth}
      \scriptsize
      \begin{tabularx}{\linewidth}{@{}l L@{}}
        \toprule
        \cabecalho{Serviço} & \cabecalho{Para quê} \\
        \midrule
        \textbf{Authentication} & login (e-mail, Google, telefone, anônimo\ldots) \\
        \textbf{Cloud Firestore} & banco NoSQL de documentos, tempo real \\
        \textbf{Cloud Storage} & fotos, PDFs, arquivos grandes \\
        \textbf{Cloud Messaging (FCM)} & notificações push \\
        \textbf{Cloud Functions} & código no servidor disparado por eventos \\
        \textbf{Crashlytics} & relatórios de falhas em produção \\
        \textbf{Analytics} & eventos e funis de uso \\
        \textbf{Remote Config} & mudar parâmetros sem publicar versão \\
        \textbf{App Check} & garante que só o seu app acessa o backend \\
        \textbf{Firebase AI Logic} & acesso a modelos Gemini a partir do app (Módulo 14) \\
        \bottomrule
      \end{tabularx}
    \end{column}
    \begin{column}{0.42\textwidth}
      \begin{conceito}[Planos]
        \textbf{Spark}: gratuito, sem cartão; ao atingir a cota, o serviço para até o próximo período.\\
        \textbf{Blaze}: pague pelo uso, com as mesmas cotas gratuitas.
      \end{conceito}
      \begin{atencao}[Cloud Storage exige Blaze]
        Desde \textbf{03/02/2026}, o Cloud Storage para Firebase só funciona em projetos no plano Blaze (a cota gratuita continua valendo). Ao ativar o Blaze, configure um \textbf{alerta de orçamento}.
      \end{atencao}
    \end{column}
  \end{columns}
\end{frame}

% =============================================================================
\section{Configuração}

\begin{frame}[fragile]{Adicionando o Firebase ao app}
  \begin{columns}[T]
    \begin{column}{0.46\textwidth}
      \begin{enumerate}
        \item Crie um projeto no console do Firebase.
        \item Adicione um app Android com o \cod{applicationId} do app.
        \item Informe o SHA-1 de debug (\cod{./gradlew signingReport}), exigido pelo login Google.
        \item Baixe o \arq{google-services.json} para \arq{app/}.
        \item Adicione o plugin e as bibliotecas.
      \end{enumerate}
      \begin{dica}
        \textbf{Tools $\rightarrow$ Firebase} no Android Studio automatiza esses passos.
      \end{dica}
    \end{column}
    \begin{column}{0.52\textwidth}
\begin{codigo}[title=app/build.gradle.kts]{Kotlin}
plugins {
  alias(libs.plugins.google.services)
  alias(libs.plugins.crashlytics)
}
dependencies {
  // o BoM alinha as versões de tudo
  implementation(platform(libs.firebase.bom))
  implementation(libs.firebase.auth)
  implementation(libs.firebase.firestore)
  implementation(libs.firebase.storage)
  implementation(libs.firebase.messaging)
  implementation(libs.firebase.crashlytics)
}
\end{codigo}
      \begin{atencao}
        Os antigos módulos \cod{-ktx} foram descontinuados: as APIs Kotlin (\cod{Firebase.auth}, \cod{snapshots()}\ldots) agora estão nos módulos principais.
      \end{atencao}
    \end{column}
  \end{columns}
\end{frame}

% =============================================================================
\section{Authentication}

\begin{frame}[fragile]{Login com e-mail e senha}
  \begin{columns}[T]
    \begin{column}{0.56\textwidth}
\begin{codigo}[title=AuthRepository.kt]{Kotlin}
class AuthRepository @Inject constructor() {
  private val auth = Firebase.auth

  // usuário atual como Flow (null = deslogado)
  val usuario: Flow<FirebaseUser?> = callbackFlow {
    val l = FirebaseAuth.AuthStateListener {
      trySend(it.currentUser)
    }
    auth.addAuthStateListener(l)
    awaitClose { auth.removeAuthStateListener(l) }
  }

  suspend fun entrar(email: String, senha: String) = runCatching {
    auth.signInWithEmailAndPassword(email, senha).await()
  }
  suspend fun cadastrar(email: String, senha: String) = runCatching {
    auth.createUserWithEmailAndPassword(email, senha).await()
  }
  fun sair() = auth.signOut()
}
\end{codigo}
    \end{column}
    \begin{column}{0.42\textwidth}
      \begin{conceito}[\cod{await()}]
        As APIs do Firebase devolvem \cod{Task}. A extensão \cod{await()} (de \cod{kotlinx-coroutines-play-services}) as transforma em chamadas \cod{suspend}.
      \end{conceito}
      \begin{dica}[Navegação condicionada]
        Observe \cod{usuario} no ViewModel principal: \cod{null} leva ao login; caso contrário, ao painel. A sessão persiste entre aberturas do app.
      \end{dica}
      \begin{atencao}
        O app nunca guarda a senha: o Firebase cuida disso e da recuperação.
      \end{atencao}
    \end{column}
  \end{columns}
\end{frame}

\begin{frame}{Login com Google via Credential Manager}
  \begin{center}
  \begin{tikzpicture}[node distance=6mm]
    \node[caixa,minimum width=2.4cm] (b) {Botão\\ ``Entrar com Google''};
    \node[caixaambar,minimum width=2.4cm,right=of b] (cm) {Credential Manager\\ \tiny seletor de contas};
    \node[caixa,minimum width=2.4cm,right=of cm] (tok) {ID token\\ \tiny do Google};
    \node[caixaescura,minimum width=2.4cm,right=of tok] (fb) {\cod{signInWithCredential}\\ \tiny Firebase Auth};
    \draw[seta] (b) -- (cm); \draw[seta] (cm) -- (tok); \draw[seta] (tok) -- (fb);
  \end{tikzpicture}
  \end{center}
  \begin{columns}[T]
    \begin{column}{0.48\textwidth}
      \begin{conceito}[Credential Manager]
        API do Android que unifica senhas, \emph{passkeys} e login federado (Google) numa única folha de seleção. Substitui o antigo \emph{Google Sign-In for Android}.
      \end{conceito}
    \end{column}
    \begin{column}{0.48\textwidth}
      \begin{dica}[Passos]
        \begin{enumerate}
          \item Ative o provedor Google no console.
          \item Pegue o \emph{Web client ID} do projeto.
          \item Peça uma credencial \cod{GetGoogleIdOption} ao \cod{CredentialManager}.
          \item Troque o ID token por uma credencial Firebase.
        \end{enumerate}
      \end{dica}
    \end{column}
  \end{columns}
  \begin{atencao}
    O SHA-1 cadastrado no console precisa corresponder à chave que assinou o app: debug, upload e a chave da Play App Signing (Módulo 12).
  \end{atencao}
\end{frame}

% =============================================================================
\section{Cloud Firestore}

\begin{frame}{Modelo de dados: coleções e documentos}
  \begin{columns}[c]
    \begin{column}{0.54\textwidth}
      \centering
      \begin{tikzpicture}[doc/.style={caixa,font=\tiny,align=left,text width=2.4cm}]
        \node[caixaescura,font=\tiny,minimum width=2.4cm] (c1) at (0,0) {coleção \cod{bancadas}};
        \node[doc] (d1) at (0.6,-1.15) {doc \cod{b3}\\ nome: ``Bancada 3''\\ sala: ``L-12'', limite: 30};
        \node[caixapetroleo,font=\tiny,minimum width=2.4cm] (c2) at (0.6,-2.3) {subcoleção \cod{leituras}};
        \node[doc] (d2) at (1.2,-3.3) {doc \cod{auto-id}\\ valor: 24.8, instante};
        \node[caixaescura,font=\tiny,minimum width=2.4cm] (c3) at (4.0,0) {coleção \cod{reservas}};
        \node[caixaambar,font=\tiny,align=left,text width=2.4cm] (d3) at (4.0,-1.25) {doc \cod{auto-id}\\ bancadaId: ``b3''\\ uid: ``x9F\ldots''\\ inicio, fim};
        \draw[seta] (c1) -- (d1); \draw[seta] (d1) -- (c2); \draw[seta] (c2) -- (d2);
        \draw[seta] (c3) -- (d3);
      \end{tikzpicture}
    \end{column}
    \begin{column}{0.44\textwidth}
      \begin{itemize}
        \item \termo{Documento}: conjunto de campos (como um JSON), até 1\,MiB.
        \item \termo{Coleção}: conjunto de documentos.
        \item Documentos podem ter \termo{subcoleções}.
        \item Não há \emph{joins}: modele pensando nas \textbf{consultas} que a tela fará.
      \end{itemize}
      \begin{dica}
        É comum \textbf{duplicar} dados (ex.: o nome da bancada dentro da reserva) para evitar leituras extras. O custo é cobrado por documento lido.
      \end{dica}
    \end{column}
  \end{columns}
\end{frame}

\begin{frame}[fragile]{Gravando e lendo em tempo real}
  \begin{columns}[T]
    \begin{column}{0.56\textwidth}
\begin{codigo}[title=ReservaRepository.kt]{Kotlin}
data class ReservaDoc(        // construtor vazio p/ Firestore
  val bancadaId: String = "", val uid: String = "",
  val inicio: Timestamp? = null, val fim: Timestamp? = null)

class ReservaRepository @Inject constructor() {
  private val col = Firebase.firestore.collection("reservas")

  suspend fun criar(r: ReservaDoc): String =
    col.add(r).await().id

  // reemite a cada mudança no servidor ou no cache
  fun daBancada(id: String): Flow<List<ReservaDoc>> = col
    .whereEqualTo("bancadaId", id)
    .whereGreaterThan("fim", Timestamp.now())
    .orderBy("fim")
    .snapshots()
    .map { it.toObjects<ReservaDoc>() }

  suspend fun cancelar(docId: String) {
    col.document(docId).delete().await()
  }
}
\end{codigo}
    \end{column}
    \begin{column}{0.42\textwidth}
      \begin{conceito}[\cod{snapshots()}]
        Transforma o \emph{listener} em tempo real do Firestore em \cod{Flow}. Quando outra pessoa cria uma reserva, todas as telas abertas atualizam.
      \end{conceito}
      \begin{conceito}[Offline]
        No Android o cache local vem ativado: leituras funcionam sem rede e escritas são enviadas quando a conexão volta.
      \end{conceito}
      \begin{atencao}[Índices]
        Filtros em campos diferentes exigem índice composto; o erro no Logcat traz o link para criá-lo.
      \end{atencao}
    \end{column}
  \end{columns}
\end{frame}

\begin{frame}[fragile]{Transações: evitando reservas duplicadas}
  \begin{columns}[T]
    \begin{column}{0.56\textwidth}
\begin{codigo}[title=Reserva atômica]{Kotlin}
suspend fun reservarHorario(b: String, slot: String,
                            uid: String) {
  val db = Firebase.firestore
  val ref = db.collection("bancadas").document(b)
    .collection("slots").document(slot)  // ex.: 2026-11-03T14

  db.runTransaction { tx ->
    val atual = tx.get(ref)
    if (atual.exists()) {
      throw FirebaseFirestoreException("Horário ocupado",
        FirebaseFirestoreException.Code.ABORTED)
    }
    tx.set(ref, mapOf("uid" to uid,
      "criadoEm" to FieldValue.serverTimestamp()))
  }.await()
}
\end{codigo}
    \end{column}
    \begin{column}{0.42\textwidth}
      \begin{conceito}[Transação]
        Lê e grava como uma operação única. Se outro cliente alterar o documento no meio, o Firestore repete a transação, e duas pessoas nunca ficam com o mesmo horário.
      \end{conceito}
      \begin{dica}[Modelagem]
        Usar o horário como \textbf{ID do documento} torna o conflito impossível por construção: só pode existir um \cod{slots/2026-11-03T14}.
      \end{dica}
      \begin{atencao}
        Transações precisam de conexão com o servidor; não funcionam offline.
      \end{atencao}
    \end{column}
  \end{columns}
\end{frame}

% =============================================================================
\section{Regras de segurança}

\begin{frame}[fragile]{Regras do Firestore}
  \begin{columns}[T]
    \begin{column}{0.56\textwidth}
\begin{codigo}[title=firestore.rules]{JSON}
rules_version = '2';
service cloud.firestore {
  match /databases/{database}/documents {

    match /bancadas/{id} {
      allow read: if request.auth != null;
      allow write: if request.auth.token.admin == true;
    }

    match /reservas/{rid} {
      allow read: if request.auth != null;
      allow create: if request.auth != null
        && request.resource.data.uid == request.auth.uid
        && request.resource.data.fim
             > request.resource.data.inicio;
      allow delete: if resource.data.uid == request.auth.uid;
    }
  }
}
\end{codigo}
    \end{column}
    \begin{column}{0.42\textwidth}
      \begin{erro}[O erro mais grave]
        Publicar com \cod{allow read, write: if true;} (``modo de teste''). Qualquer pessoa na internet lê e apaga todo o banco.
      \end{erro}
      \begin{conceito}[Como pensar]
        O app \textbf{não é confiável}: alguém pode chamar a API do Firebase diretamente. As regras são a única barreira real e rodam no servidor.
      \end{conceito}
      \begin{dica}[Emulator Suite]
        \cod{firebase emulators:start} roda Auth, Firestore e Storage localmente, sem custo e sem dados reais.
      \end{dica}
    \end{column}
  \end{columns}
\end{frame}

% =============================================================================
\section{Storage, mensagens e monitoramento}

\begin{frame}[fragile]{Cloud Storage: fotos dos equipamentos}
  \begin{columns}[T]
    \begin{column}{0.54\textwidth}
\begin{codigo}[title=FotoRepository.kt]{Kotlin}
suspend fun enviarFoto(patrimonio: String, uri: Uri): String {
  val ref = Firebase.storage.reference
    .child("equipamentos/$patrimonio.jpg")
  ref.putFile(uri).await()                // upload
  return ref.downloadUrl.await().toString()
}

// No ViewModel: grava a URL no documento do equipamento
val url = fotos.enviarFoto("IFMA-0042", uriEscolhida)
equipamentos.atualizarFoto("IFMA-0042", url)
\end{codigo}
      \begin{dica}
        Comprima a imagem antes de enviar (ex.: 1280\,px, JPEG 80\,\%). Fotos de câmera têm vários MB.
      \end{dica}
    \end{column}
    \begin{column}{0.44\textwidth}
\begin{codigo}[title=storage.rules]{JSON}
match /equipamentos/{arquivo} {
  allow read: if request.auth != null;
  allow write: if request.auth != null
    && request.resource.size < 5 * 1024 * 1024
    && request.resource.contentType
         .matches('image/.*');
}
\end{codigo}
      \begin{conceito}
        Arquivos ficam no Storage; o Firestore guarda apenas a URL e os metadados.
      \end{conceito}
    \end{column}
  \end{columns}
\end{frame}

\begin{frame}[fragile]{Cloud Messaging: notificações push}
  \begin{columns}[T]
    \begin{column}{0.5\textwidth}
      \centering
      \begin{tikzpicture}[node distance=5mm]
        \node[caixaescura,minimum width=2.6cm] (srv) {Servidor / Function\\ \tiny ``Bancada 3 acima de 35\,°C''};
        \node[caixaambarforte,minimum width=2.6cm,below=of srv] (fcm) {FCM};
        \node[below=of fcm] (cel) {\celular[0.32]{}};
        \draw[seta] (srv) -- node[rotulo,right]{token do aparelho} (fcm);
        \draw[seta] (fcm) -- (cel);
      \end{tikzpicture}
\begin{codigo}[title=AndroidManifest.xml]{XML}
<uses-permission android:name=
  "android.permission.POST_NOTIFICATIONS"/>
\end{codigo}
    \end{column}
    \begin{column}{0.48\textwidth}
\begin{codigo}[title=MensagensService.kt]{Kotlin}
class MensagensService : FirebaseMessagingService() {
  override fun onNewToken(token: String) {
    // envie ao seu backend / Firestore
    salvarTokenDoUsuario(token)
  }
  override fun onMessageReceived(msg: RemoteMessage) {
    val titulo = msg.notification?.title ?: "LabMonitor"
    mostrarNotificacao(titulo, msg.data["texto"])
  }
}
\end{codigo}
      \begin{atencao}[Android 13+]
        Notificações exigem a permissão em tempo de execução \cod{POST\_NOTIFICATIONS} (Módulo 10). Registre o \emph{service} no manifesto.
      \end{atencao}
    \end{column}
  \end{columns}
\end{frame}

\begin{frame}{Crashlytics, Analytics e alternativas ao Firebase}
  \begin{columns}[T]
    \begin{column}{0.48\textwidth}
      \begin{conceito}[Crashlytics]
        Agrupa as falhas de produção por causa, com pilha de chamadas, modelo do aparelho e versão do Android. Use \cod{Firebase.crashlytics.log("...")} e \cod{setCustomKey} para dar contexto.
      \end{conceito}
      \begin{conceito}[Analytics]
        \cod{logEvent("reserva\_criada")} registra eventos e mede o uso real das funcionalidades.
      \end{conceito}
      \begin{atencao}[LGPD]
        Informe na política de privacidade quais dados são coletados e por quê (Módulo 12).
      \end{atencao}
    \end{column}
    \begin{column}{0.5\textwidth}
      \scriptsize
      \begin{tabularx}{\linewidth}{@{}l L@{}}
        \toprule
        \cabecalho{Alternativa} & \cabecalho{Características} \\
        \midrule
        \textbf{Supabase} & PostgreSQL (SQL e relações), auth, storage, tempo real; código aberto \\
        \textbf{Appwrite} & código aberto, pode ser hospedado no próprio servidor \\
        \textbf{Backend próprio} & Ktor, Spring, FastAPI ou Node com banco relacional; controle total, mais trabalho \\
        \bottomrule
      \end{tabularx}
      \begin{dica}[Mantenha a porta aberta]
        Com a arquitetura do Módulo 06, o Firebase fica atrás de interfaces de repositório. Trocar de backend afeta só a camada de dados.
      \end{dica}
    \end{column}
  \end{columns}
\end{frame}

% =============================================================================
\section{Fechamento}

\begin{frame}{Resumo}
  \begin{resumo}
    \begin{itemize}
      \item Firebase entrega autenticação, banco, arquivos, push e monitoramento sem servidor próprio.
      \item Configure com \arq{google-services.json}, o plugin e o BoM; use \cod{await()} e \cod{snapshots()} com coroutines.
      \item Firestore: coleções e documentos, consultas com índices, cache offline e transações.
      \item As regras de segurança são a proteção real dos dados; teste-as no Emulator Suite.
      \item Storage guarda arquivos (exige plano Blaze); FCM entrega notificações push.
      \item Isole o Firebase atrás de repositórios para poder trocar de backend.
    \end{itemize}
  \end{resumo}
\end{frame}

\begin{frame}{Exercícios}
  \begin{exercicio}[Exercício 1 --- Autenticação]
    Implemente as telas de login, cadastro e ``esqueci a senha'' (\cod{sendPasswordResetEmail}) com validação e mensagens de erro em português.
  \end{exercicio}
  \begin{exercicio}[Exercício 2 --- Regras]
    Escreva regras que permitam a qualquer usuário logado ler as bancadas, mas só o dono editar ou cancelar a própria reserva. Teste no emulador com dois usuários.
  \end{exercicio}
  \begin{exercicio}[Exercício 3 --- Tempo real]
    Abra o app em dois aparelhos com usuários diferentes. Crie uma reserva em um e observe o outro. Desligue a rede de um deles e explique o comportamento.
  \end{exercicio}
  \begin{projeto}[LabMonitor --- Etapa 9]
    Adicione login, reservas de bancada no Firestore com transação anti-conflito e regras de segurança. Integre o Crashlytics.
  \end{projeto}
\end{frame}

\begin{frame}{Para aprofundar}
  \small
  \begin{itemize}
    \item Firebase para Android: \url{https://firebase.google.com/docs/android/setup}
    \item Authentication: \url{https://firebase.google.com/docs/auth/android/start}
    \item Credential Manager: \url{https://developer.android.com/identity/sign-in/credential-manager}
    \item Cloud Firestore: \url{https://firebase.google.com/docs/firestore}
    \item Regras de segurança: \url{https://firebase.google.com/docs/rules}
    \item Emulator Suite: \url{https://firebase.google.com/docs/emulator-suite}
    \item Preços: \url{https://firebase.google.com/pricing}
  \end{itemize}
\end{frame}

\encerramento{Módulo 10 --- Recursos do dispositivo e IoT: permissões, câmera, localização, WorkManager, sensores, BLE e MQTT com ESP32.}

\end{document}
